\chapter{Introduccion}

\subsection{Contexto}

El sistema electrico chileno (SEN) alimenta a mas de 18 millones de habitantes
a traves de una red de transmision troncal y mas de 600 subestaciones de
distribucion que sirven a las principales concesionarias del pais:
Chilquinta (Valparaiso), CGE (zona central-norte), Saesa (sur), Enel Distribucion
(RM) y otras. La continuidad del servicio se mide por los indicadores SAIDI
(duracion media de interrupcion) y SAIFI (frecuencia media), que son
fiscalizados por la SEC y expuestos anualmente al SERNAC.

El frente de mal tiempo de julio de 2026 en Valparaiso, con rafagas de viento
superiores a 80 km/h y lluvias intensas que provocaron la caida de lineas de
23~kV y la salida no programada de subestaciones, motivo cargos directos de la
SEC contra Chilquinta y un procedimiento del SERNAC por reposicion tardia en
cuadrillas. La fiscalizacion post-evento y la trazabilidad de las decisiones
de despacho son ahora una prioridad operativa y reputacional.

Las subestaciones modernas estan equipadas con:
\begin{itemize}[leftmargin=*,itemsep=2pt]
  \item \textbf{SCADA}: telemetría a 1-5~s de tension, corriente, potencia,
        posicion de interruptores, alarma de protecciones.
  \item \textbf{PMU} (\emph{Phasor Measurement Units}): sincrofasores a
        25-50~fps con precision sincronizada por GPS, dan una vision
        electromecanica fina del sistema.
  \item Datos de calidad (THD, flicker, desbalance) y temperatura de
        devanados de transformadores.
\end{itemize}

Estos datos son \emph{ideales} para entrenar modelos de deteccion temprana
de anomalias, pero su uso en la practica enfrenta tres obstaculos que
definen el problema de este trabajo:

\begin{enumerate}[leftmargin=*,itemsep=4pt]
  \item \textbf{Etiquetas casi inexistentes}. Las fallas son raras
        ($\lambda \approx 0.5$/dia por subestacion segun statistic de SEC).
        Aunque la falla ocurre, muchas veces la senal previa es ambigua
        y los operadores la marcan \emph{a posteriori} solo si la
        proteccion opero. Esto inutiliza los enfoques de clasificacion
        supervisada convencionales.

  \item \textbf{Datos que no pueden salir del recinto}. Cada
        distribuidora considera sus datos como activo comercial y la
        comparticion SER/NAL es voluntaria. Ademas, la normativa
        NERC CIP y los estandares CNE chilenos exigen segregacion
        estricta entre redes OT (operativas) e IT (corporativas).
        Esto descarta el entrenamiento centralizado en la nube.

  \item \textbf{Eventos climaticos extremos} cada vez mas frecuentes.
        Frentes de viento, olas de calor y tormentas obligan a las
        distribuidoras a \emph{priorizar} cuadrillas de reposicion.
        Un sistema que solo alerte de fallas individuales es
        insuficiente; hace falta un ranking de riesgo de
        reposicion lenta por subestacion.
\end{enumerate}

\subsection{Propuesta}

Este trabajo responde con un pipeline de cuatro componentes que se
describen en los capitulos siguientes:

\begin{enumerate}[leftmargin=*,itemsep=2pt]
  \item \textbf{Gemelo digital ligero}: modelo fisico-estocastico de la
        subestacion (estado estacionario + dinamica termica y de
        regulador) que predice el comportamiento esperado. El gemelo
        tambien genera datos sinteticos etiquetados para entrenar el
        detector sin depender de fallas reales etiquetadas.

  \item \textbf{Detector LSTM autoencoder + conformal prediction}:
        aprende la dinamica normal del vector de desviaciones observado
        versus esperado. La salida es un reconstruction error; sobre el
        se aplica conformal prediction para fijar el umbral de
        deteccion con una \emph{tasa de falsa alarma garantizada}
        $\alpha = 0.05$.

  \item \textbf{Entrenamiento federado (FedAvg)} entre cinco
        subestaciones virtuales. Cada SE entrena un autoencoder local;
        un orquestador promedia los pesos. Los datos crudos nunca
        salen del recinto, cumpliendo con la normativa y preservando
        la propiedad.

  \item \textbf{Priorizacion post-evento}: tras un evento climatico
        extremo (frente de viento, ola de calor), el sistema entrega un
        orden de despacho de cuadrillas por subestacion, ponderando
        detecciones recientes, estado termico y criticidad geografica.
\end{enumerate}

\subsection{Aportes}

Los aportes originales son:

\begin{itemize}[leftmargin=*,itemsep=2pt]
  \item Demostracion de que un \emph{gemelo digital simplificado} de
        subestacion produce residuos con suficiente senal/ruido para
        entrenar detectores no supervisados en el dominio de distribucion
        chilena.

  \item Validacion experimental (sobre simulacion calibrada con datos
        CEN publicos) de que un \emph{autoencoder LSTM} entrenado
        sobre los residuos del gemelo alcanza un FPR empirico de
        $4.86\%$ contra el objetivo conformal del $5\%$, con recall
        de $76$-$80\%$ y latencia mediana de deteccion nula,
        preservando la confianza del operador.

  \item Cuantificacion del \emph{costo de privacidad}: en el
        prototipo no hay costo medible -- el modelo federado iguala
        o supera al centralizado -- mientras que los datos nunca
        abandonan el recinto.

  \item Una \emph{arquitectura de producto} lista para empaquetar como
        SaaS para distribuidoras y transmisoras chilenas, con un caso
        de uso directo para el cumplimiento regulatorio post-Frente
        Valparaiso 2026.
\end{itemize}

\subsection{Organizacion del documento}

\begin{itemize}[leftmargin=*,itemsep=2pt]
  \item \textbf{Capitulo 2}: marco teorico (gemelos digitales en
        sistemas electricos de potencia, conformal prediction,
        aprendizaje federado, regulacion electrica chilena).
  \item \textbf{Capitulo 3}: datos (CEN, gemelo digital, catalogo
        de fallas).
  \item \textbf{Capitulo 4}: metodologia detallada del pipeline.
  \item \textbf{Capitulo 5}: resultados experimentales.
  \item \textbf{Capitulo 6}: discusion.
  \item \textbf{Capitulo 7}: conclusiones y trabajo futuro.
\end{itemize}

El notebook y el dashboard asociados a este trabajo se entregan como
material complementario.